%% 导入宏包
\usepackage[a4paper, scale=0.8]{geometry} %纸张大小和页边距
\usepackage{fancyhdr} %页眉和页脚
\usepackage{graphicx} %导入图片，graphicx是graphics的拓展，graphics自动加载
\usepackage{amsmath, amssymb, amsthm} %amsmath提供数学公式的优化以及部分功能，amssymb定义了新的数学符号，amsthm提供了数学环境，amssymb包含了amsfonts
\usepackage{mathtools} %mathtools包含了amsmath，加在mathtools可以不加载amsmath
\usepackage[mathrm=sym,warnings-off={mathtools-colon,mathtools-overbracket}]{unicode-math} %提供数学字体
\usepackage{enumitem} %调整列表的宏包
\usepackage{physics} %物理学公式常用宏包，可能与sinuitx冲突(\qty命令)
\usepackage[version=3]{mhchem} %化学宏包，主要用来写同位素核素等符号，version=3是为了不和lyx冲突
\usepackage{framed} %加边框
\usepackage{float} %让浮动体强制固定在此处
\usepackage{subfigure} %子图宏包
\usepackage{listings} %代码环境
\usepackage{import} %主要为配合incfig命令使用
\hypersetup{unicode=true} %超链接功能
\renewcommand*{\backref}[1]{}
\renewcommand*{\backreftwosep}{ 与~}
\renewcommand*{\backreflastsep}{，以及~}
\renewcommand*{\backrefsep}{，}
\renewcommand*{\backrefalt}[4]{%
  \ifcase #1 %
  \or
    \textit{（见第~#2~页）}%
  \else
    \textit{（见第~#2~页）}%
  \fi
}
\usepackage{tensor} %张量功能
\usepackage[export]{adjustbox} %为了在公式中使用图片
\usepackage{Titlepage} %绘制封面
\usepackage{simpler-wick} %wick缩并功能，调用了tikz宏包
\usepackage{tikz} %绘图宏包，由于titlepage中导入且可选项不同所以不加
\usepackage{perpage} %the perpage package
\MakePerPage{footnote}
\usepackage{multicol} %多栏排版（用于主题索引）
\usepackage{longtable} %长表格（用于专有术语对照表）
\usepackage{booktabs} %三线表优化
\usepackage{etoolbox} %环境钩子宏包
\AtBeginEnvironment{table}{\centering} %默认所有表格浮动体自动居中
\AtBeginEnvironment{table*}{\centering}

% 索引与术语表排版宏
\newcommand{\indexentry}[3]{%
  \noindent\hangindent=1.5em\hangafter=1\textbf{#1}\ifstrempty{#2}{}{（#2）}\dotfill #3\par
}
\newcommand{\indexparent}[2]{%
  \noindent\hangindent=1.5em\hangafter=1\textbf{#1}\ifstrempty{#2}{}{（#2）}\par
}
\newcommand{\indexsubentry}[3]{%
  \noindent\hangindent=2.5em\hangafter=1\hspace*{1.2em}\textendash~#1\ifstrempty{#2}{}{（#2）}\dotfill #3\par
}
\newcommand{\indexentrysee}[3]{%
  \noindent\hangindent=1.5em\hangafter=1\textbf{#1}\ifstrempty{#2}{}{（#2）}\dotfill\textit{#3}\par
}

%% 内容格式设置
%利用Unicode-math的字体解决方案，来源：许兆晖
\setmainfont{texgyrepagella-regular.otf}[
    BoldFont = texgyrepagella-bold.otf,
    ItalicFont = texgyrepagella-italic.otf,
    BoldItalicFont = texgyrepagella-bolditalic.otf
]
\setsansfont{texgyreheros-regular.otf}[
    BoldFont = texgyreheros-bold.otf,
    ItalicFont = texgyreheros-italic.otf,
    BoldItalicFont = texgyreheros-bolditalic.otf
]
\setmonofont{texgyrecursor-regular.otf}[
    BoldFont = texgyrecursor-bold.otf,
    ItalicFont = texgyrecursor-italic.otf,
    BoldItalicFont = texgyrecursor-bolditalic.otf
]
\setmathfont{texgyrepagella-math.otf}[AutoFakeBold = 2.5]
\setmathfont{XITSMath-Regular.otf}[
    BoldFont = XITSMath-Bold.otf,
    range = {scr, bfscr},
]
\setmathfont{XITSMath-Regular.otf}[
    BoldFont = XITSMath-Bold.otf,
    FakeSlant = -.25,
    range = \ell,
]
% Prefer Cambria when available; use bundled Pagella Math on other systems.
\IfFontExistsTF{Cambria Math}{
\setmathfont{Cambria Math}[
    range = {\setminus},
]
\setmathfont{Cambria Math}[
    range = {cal,bfcal},
    Scale = 1.05,
]
}{
\setmathfont{texgyrepagella-math.otf}[
    range = {\setminus},
]
\setmathfont{texgyrepagella-math.otf}[
    range = {cal,bfcal},
    Scale = 1.05,
]
}

% 设置数学字体的斜粗体
\renewcommand{\boldsymbol}[1]{
    \begingroup
        \let\symup\symbfup\let\mathrm\mathbfup
        \let\symit\symbfit\let\mathit\mathbfit
        \let\symbf\symbfsf\let\mathsf\mathbfsf
        \symbfit{#1}
    \endgroup
}

% 设置页眉与页脚，需要fancyhdr宏包
\pagestyle{fancy}
\lhead{技术物理系}
\chead{}
\rhead{}
\lfoot{}
\cfoot{}
\rfoot{\thepage}
\setlength{\headheight}{15pt} %设置页眉到页面顶部的长度，可能需要geometry宏包
% 设置列表样式，需要enumitem宏包
\setitemize[1]{itemsep=0pt,partopsep=0pt,parsep=\parskip,topsep=5pt}
% 设置为无首行缩进
\setlength{\parindent}{0pt}
% 设置图片路径
\graphicspath{{pictures/}}
%% 非必要的设置
% 设置代码样式，需要listings宏包
\lstset{
    language = Python,
    numbers=left, 
    numberstyle= \small, 
    keywordstyle= \color{green},
    commentstyle= \color{brown}, 
    frame=shadowbox, % 阴影效果
    rulesepcolor= \color{green!20!} ,
    escapeinside=``, % 英文分号中可写入中文
    xleftmargin=2em,xrightmargin=2em, aboveskip=1em,
    framexleftmargin=2em,
    framexrightmargin = 1em,
    breaklines %自动换行
}
\definecolor{green}{RGB}{57,174,169}
\definecolor{brown}{RGB}{202,149,92}

%% 自定义的数学有关命令
\AtBeginDocument{
    \renewcommand{\Re}{\mathop{\text{Re}}}
    \renewcommand{\Im}{\mathop{\text{Im}}}
}
% 交叉引用的设置
\AtBeginDocument{
\let\ref\autoref
\renewcommand\equationautorefname{式}
\renewcommand\figureautorefname{图}
\renewcommand\tableautorefname{表}
\renewcommand\sectionautorefname{节}
\renewcommand\subsectionautorefname{小节}
}

%% 自定义超大括号
\makeatletter
\newcommand{\biggg}{\bBigg@{3}}
\def\bigggl{\mathopen\biggg}
\def\bigggm{\mathrel\biggg}
\def\bigggr{\mathclose\biggg}
\newcommand{\Biggg}{\bBigg@{3.5}}
\def\Bigggl{\mathopen\Biggg}
\def\Bigggm{\mathrel\Biggg}
\def\Bigggr{\mathclose\Biggg}
\newcommand{\bigggg}{\bBigg@{4}}
\def\biggggl{\mathopen\bigggg}
\def\biggggm{\mathrel\bigggg}
\def\biggggr{\mathclose\bigggg}
\newcommand{\Bigggg}{\bBigg@{4.5}}
\def\Biggggl{\mathopen\Bigggg}
\def\Biggggm{\mathrel\Bigggg}
\def\Biggggr{\mathclose\Bigggg}
\makeatother